\section{总结与延伸}

\subsection{核心要点}

\begin{enumerate}
    \item 复合 AI 系统通过组合多个组件突破单次 LLM 调用的能力限制
    \item 手工 prompting 在复合系统中脆弱、不可组合、难以优化
    \item DSPy 提供声明式编程范式：定义结构 + 自动编译 prompt
    \item Signature、Module、Optimizer、Metric 是四个核心抽象
    \item 类比编译器：高级语言（DSPy 程序）→ 底层实现（优化后的 prompt）
\end{enumerate}

\subsection{拓展阅读}

\begin{itemize}
    \item Khattab et al., ``DSPy: Compiling Declarative Language Model Calls into Self-Improving Pipelines,'' ICLR 2024.
    \item Khattab et al., ``Demonstrate-Search-Predict: Composing Retrieval and Language Models for Knowledge-Intensive NLP,'' 2023.
\end{itemize}

\subsection{全局总结表}

\begin{table}[H]
\centering
\begin{tabular}{p{0.22\textwidth}p{0.25\textwidth}p{0.25\textwidth}}
\toprule
主题 & 能力提升 & 工程风险 \\
\midrule
Compound Systems & 管道化多模块 & 组合漏洞成为链式故障 \\
DSPy 编译 & 声明式 prompt 生成 & 搜索空间大，需可追踪结果 \\
Evidence/Monitoring & 让每次调用可观察 & 监控缺位会放大退化 \\
\bottomrule
\end{tabular}
\caption{Lecture 08 的价值/风险对照}
\end{table}

\subsection{未来行动}

\begin{itemize}
    \item 为每个 optimizer 版本建立 baseline 与 rollout log
    \item 把 Evidence Matrix 绑定报警系统，自动捕捉 prompt drift
    \item 保持 Action Timeline up-to-date，以便快速定位责任人
\end{itemize}


\end{document}
